\section{Simultaneous private top-rank depths in all three actual arms}
\label{td-section}

The sufficient signed-compensation classes do not have automatic
membership. This section constructs actual primitive unramified roots
whose three quotient arms simultaneously have prescribed private deep
primes above the moving cutoff. It also accounts for the height paid by
the construction. Neither the net-credit criterion nor the full signed
tail is refuted.

\begin{theorem}[Three prescribed actual first depths]\label{td-three-depths}
Let $n\ge5$ be an integer coprime to six, and let $s_1,s_2,s_3\ge4$
be integers. There are distinct primes $q_i\equiv1\pmod{6n}$ and
infinitely many positive integers $a\equiv0\pmod3$ such that
$w=a+\zeta$ is primitive and unramified and has the following properties.
Form the normalized power
\[
 z_n=A+B\zeta=
 \begin{cases}w^n,&n\equiv1\pmod6,\\
 \bar w^{\,n},&n\equiv5\pmod6,
 \end{cases}
 \qquad
 D_1=A/a,\quad D_2=B,\quad D_3=(A+B)/(a+1).
\]
These are the actual integer quotients of Lemma~\ref{sac-quotients}.
For $i\ne j$,
\begin{equation}\label{td-depths}
 v_{q_i}(D_i)=s_i,\qquad v_{q_i}(D_j)=0,\qquad
 q_i\nmid a(a+1)N(w),\qquad
 \operatorname{ord}\bigl((w/\bar w)^3\bmod q_i\bigr)=n.
\end{equation}
In particular $v_{q_i}(T_n)=s_i$ for $T_n=|P(w^n)|$, and these
are genuinely new first depths at rank $n$.
\end{theorem}

\begin{proof}
We first give an elementary supply of distinct primes. For any integer
$N\ge3$ and finite set of previously chosen primes, put $X$ equal to
$N$ times their product. The integer $\Phi_N(X)$ is greater than one:
pairing primitive roots with their conjugates gives positive factors
at least $(X-1)^2>1$. Its constant term is one, so any prime divisor
$q$ is coprime to $X$, and hence to $N$ and all the previous primes.
The polynomial $Z^N-1$ is separable over $\mathbb F_q$. Its cyclotomic
factors are therefore pairwise coprime after reduction. A root of
$\Phi_N$ cannot also be a root of $Z^d-1$ at a proper divisor $d$ of
$N$. Thus $X\bmod q$ has order exactly $N$, and $q\equiv1\pmod N$.
Repeat this three times with $N=6n$ to obtain the required distinct
$q_i$ and elements $t_i$ of order $6n$. No prime-density theorem is used.

Fix such a pair $q,t$ and put $\theta=t^n$ in $\mathbb F_q$. It is
a primitive sixth root, with $\theta^2-\theta+1=0$ and
$\bar\theta=1-\theta=\theta^{-1}$. On the norm-unit locus define
\[
 \alpha(X)=\frac{X+\theta}{X+\bar\theta}.
\]
For $\alpha_0\ne1$ the inverse coordinate is
\begin{equation}\label{td-inverse}
 X_0=\frac{\alpha_0\bar\theta-\theta}{1-\alpha_0}.
\end{equation}
In particular
\[
 X_0+\bar\theta=\frac{\bar\theta-\theta}{1-\alpha_0},
 \qquad X_0+\theta=
   \alpha_0\frac{\bar\theta-\theta}{1-\alpha_0}.
\]
Choose $\alpha_0$ successively as $t^2,t^6,t^4$. Their $n$th powers
are $\theta^2,1,\theta^4$, respectively, and their cubes each have
order exactly $n$. Their orders exceed three, so $X_0$ is neither
zero nor minus one. The two displayed expressions are nonzero,
so the norm and all input arms are units at $X_0$.

Write the raw power as
$(X+\zeta)^n=A_n(X)+B_n(X)\zeta$. Its ratio to its conjugate is
$\alpha(X)^n$. Cross multiplication on the norm-unit locus gives
\[
 \begin{array}{c|c}
 \alpha(X)^n&\text{vanishing raw arm}\\ \hline
 \theta^2&A_n(X)\\
 1&B_n(X)\\
 \theta^4&A_n(X)+B_n(X).
 \end{array}
\]
Two arms cannot vanish together there. Each difference of ratios is
a nonzero constant times the corresponding arm divided by the unit
$(X+\bar\theta)^n$. Moreover
\[
 \alpha'(X)=\frac{\bar\theta-\theta}{(X+\bar\theta)^2}\ne0.
\]
Since $q>n$, the indicated raw arm has a simple zero at $X_0$.
Dividing by the relevant input $X$, $1$ or $X+1$, all units there,
preserves this simple zero of the actual quotient polynomial.

For $n\equiv1\pmod6$, the choices $t^2,t^6,t^4$ target
$D_1,D_2,D_3$ in that order. For $n\equiv5\pmod6$, conjugation
sends the raw coordinates to $(A_n+B_n,-B_n)$. The correct choices
are then $t^4,t^6,t^2$. This exchange is what identifies the
labeled actual quotients in \eqref{td-depths}.

An integer polynomial with a simple root modulo $q$ has exactly one
lift to a root at the next precision: for $e\ge1$,
\[
 D_i(r+jq^e)\equiv D_i(r)+jq^eD_i'(r)\pmod{q^{e+1}}.
\]
Lift the chosen root uniquely through precision $q^{s_i}$. At the
next step choose any of the $q-1$ nonroot digits. This gives a residue
$r_i\bmod q_i^{s_i+1}$ of exact depth $s_i$. The other quotient
arms, input arms and norm remain units modulo $q_i$, and the cube
ratio keeps exact order $n$.

The Chinese remainder theorem combines the three residues with
$a\equiv0\pmod3$. Write
\begin{equation}\label{td-crt-modulus}
 M_0=3\prod_{i=1}^3q_i^{s_i+1},\qquad 0\le r<M_0
\end{equation}
for the resulting modulus and representative. Every $a=r+jM_0$,
$j\ge1$, is positive and preserves all those exact depths.
The root has coordinate gcd one, and $N(w)=a^2+a+1\equiv1\pmod3$.
It is an actual primitive unramified nonunit root. Its powers have
nonzero primitive boundary by the established oriented factorization.
Since $T_1=a(a+1)$ is coprime to each chosen prime and
$T_n=T_1|D_1D_2D_3|$, the stated full depths and first ranks follow.
\end{proof}

\begin{corollary}[Failure of automatic depth-cap membership]
\label{td-cap-boundary}
The constructed roots give a family with $n\to\infty$ and unit
residual, hence $\lambda=1/n$ and $\rho=\log(4+n)/n$.
Every $q_i>n^{1/6}$. No pair of actual quotient arms has all its
prime depths above that cutoff at most two.

More generally, if caps $h_i\in\mathbb Z_{>0}\cup\{\infty\}$
bound all those actual depths, with $1/\infty=0$, then
\[
 \sum_{i=1}^3\frac1{h_i}\le\frac34<1.
\]
Thus the union of reciprocal-cap sufficient classes does not contain
every actual root, even when a cap may be selected separately for
each root. This does not refute the quantitative net-credit criteria.
\end{corollary}
\begin{proof}
The representation uses exponent $n$ and unit residual. The selected
primes are at least $6n+1$, and every selected depth is at least four.
Every finite valid cap is consequently at least four; an infinite
cap contributes zero to the reciprocal sum. The assertions follow
directly from Theorem~\ref{td-three-depths} for each allowed $n$.
\end{proof}

\begin{proposition}[The selected depths have a visible height cost]
\label{td-selected-cost}
For the representatives $a\ge M_0$ used above, put
$c=\max(|A|,|B|,|A+B|)$ and $t_{\rm ht}=\log c$. The contribution
of the selected primes alone to the excess above depth two satisfies
\begin{equation}\label{td-cost}
 E_{\rm selected}=\sum_{i=1}^3(s_i-2)\log q_i
       <\frac{t_{\rm ht}}n.
\end{equation}
The selected positive signed depths $(s_i-3)\log q_i$ have a
smaller total. The estimate is uniform even when all prescribed
depths vary with $n$.
\end{proposition}
\begin{proof}
The actual quadratic height gives $Q^n\le c^2$, while
$Q=a^2+a+1>a^2$. Hence
\[
 t_{\rm ht}>n\log a\ge n\log M_0.
\]
From \eqref{td-crt-modulus},
$E_{\rm selected}<\log M_0$, which proves \eqref{td-cost}.
\end{proof}

\paragraph{Evidence and remaining scope.}
The root researcher and both research peers independently reviewed the
complete ordinary proof. The exact finite replay checks
$n=5,7,11,25,35$, with depths $(4,5,6)$ and two actual CRT roots
per exponent: ten actual roots and thirty selected full-depth checks.
It proves the selected primes prime by trial division, checks every
required order against all proper divisors, and exhaustively checks
each selected lifting level. All actual coordinates, norm and quotient
identities, coprimality, selected depths and the integer height-cost
certificate are recomputed exactly. One other researcher independently
reread the replay and reran its canonical check successfully.

The quotient arms are not completely factored in that replay. No
simultaneous second pure norm, full signed-tail failure or Lean proof
of this construction is asserted. The unselected prime factors may
supply substantial negative radical credit. Equation~\eqref{td-cost}
also shows that these selected deep primes alone pay only a vanishing
normalized cost. General control of the actual excess-minus-radical
quantity, including unusually short congruence directions and all
unselected factors, remains open.
